\section{An admission gate that meets R1 to R4}
\label{sec:gate}

Figure~\ref{fig:protocol} shows the protocol. Each component below answers one requirement.

\begin{figure}[t]
\centering
% F1: admission protocol (TikZ, included by paper/main.tex)
\begin{tikzpicture}[font=\small, node distance=4mm and 6mm,
  box/.style={draw, rounded corners=2pt, align=center, minimum height=9mm, inner sep=3pt, fill=#1},
  box/.default=white, ar/.style={-{Stealth[length=2mm]}, thick}]
\definecolor{okblue}{HTML}{0072B2}\definecolor{okorange}{HTML}{E69F00}\definecolor{okgreen}{HTML}{009E73}
\node[box=okblue!12] (reg) {Register\\$H(s)$, $\mathrm{id}(a)$ at $b_0$};
\node[box=okblue!12, right=of reg] (com) {Commit $c=H(\sigma_p\|\rho\|\theta)$\\at $b_1$, beacon height $b_2$};
\node[box=okblue!12, right=of com] (seed) {Seed\\$\sigma=H(\sigma_p\|\beta_{b_2})$};
\node[box=okorange!18, below=7mm of seed] (exam) {Paired exam on $G(\sigma,i)$\\arm 0: $s_0$, arm 1: $s$};
\node[box=okorange!18, left=of exam] (gate) {Betting test, futility margin;\\level $\alpha_t$ from e-LOND\\and lineage spending};
\node[box=okgreen!15, left=of gate] (cred) {Credential $\chi$: both\\Merkle roots, $\theta$, decision;\\reveal $(\sigma_p,\rho)$ at $b_3$};
\node[box=okgreen!15, below=7mm of gate] (ver) {Any verifier: Algorithm~\ref{alg:verify}\\recompute decision from transcripts};
\node[box=gray!12, below=7mm of cred] (log) {Append-only log\\+ public beacon};
\draw[ar] (reg) -- (com); \draw[ar] (com) -- (seed); \draw[ar] (seed) -- (exam);
\draw[ar] (exam) -- node[above, font=\scriptsize]{$D_i$} (gate); \draw[ar] (gate) -- node[above, font=\scriptsize]{stop} (cred);
\draw[ar] (cred) -- (log); \draw[ar] (log) -- (ver); \draw[ar, dashed] (ver.east) -| (exam.south);
\end{tikzpicture}

\caption{Admission protocol. The skill digest is logged before the proctor commits to the exam parameters $\theta$ and a secret $\sigma_p$; the variant seed mixes $\sigma_p$ with a public beacon value that does not yet exist at commit time. The exam runs paired arms on variants $G(\sigma,i)$ and stops by the registered rule. The credential binds both arms' transcript Merkle roots and the decision, and any verifier can recompute the decision (Algorithm~\ref{alg:verify}). Any append-only log with a public beacon can serve (a public chain, or a transparency log plus drand).}
\label{fig:protocol}
\end{figure}

\subsection{Hidden variants by commit and reveal (R4)}
The registry logs the skill digest $H(s)$ and agent identity at log height $b_0$. The proctor then draws a secret $\sigma_p$ and salt $\rho$ and logs $c=H(\sigma_p\,\|\,\rho\,\|\,\theta)$ at $b_1>b_0$, where $\theta$ is the registered parameter block (test, $\alpha$, margins, $r$, $N_{\max}$, bet rule, generator version, beacon height $b_2>b_1$). The exam seed is $\sigma=H(\sigma_p\,\|\,\beta_{b_2})$, where $\beta_{b_2}$ is the beacon value at $b_2$. After the exam the proctor reveals $(\sigma_p,\rho)$. Because $H(s)$ precedes $c$, the skill cannot depend on the variants; because $\sigma_p$ is secret, the agent cannot predict them; because $c$ precedes $\beta_{b_2}$, the proctor cannot choose variants after seeing the skill. Commitments without credentials stay visible, so abandoned exams cannot be hidden. A colluding proctor who leaks $\sigma_p$ is not stopped by this construction; threshold proctors or seed release inside a trusted execution environment address that case (Section~\ref{sec:discussion}).

With an unpredictable seed a skill can carry answers only by covering a large share of the generator's space, either with a lookup table or with a copy of the generator or its solver; the coverage statistic below measures the first, and Section~\ref{sec:res-leak} discusses the second. The registry estimates that share, $\lambda_{\text{cover}}$, on a Monte Carlo sample of the space and rejects skills above a threshold (Section~\ref{sec:res-leak}).

\subsection{The betting test with a futility margin (R1)}
Let $X_i=(1+D_i)/2\in[0,1]$ and $m_0=(1+\delta)/2$, so $H_0$ of Eq.~\eqref{eq:H} reads $\mathbb{E}[X_i\mid\mathcal F_{i-1}]\le m_0$ for every $i$. The wealth process of \citet{waudbysmith2024betting} is
\begin{equation}
K_n=\prod_{i=1}^n\bigl(1+\lambda_i(X_i-m_0)\bigr),\qquad
\lambda_i=\min\Bigl\{c,\ \sqrt{\tfrac{2\log(1/\alpha)}{\hat v_{i-1}\, i\log(1+i)}}\Bigr\},
\label{eq:wealth}
\end{equation}
with $\hat v_{i-1}$ the predictable plug-in variance ($\hat v_0=1/4$) and cap $c\in\{1/2,3/4\}$. Here $\alpha$ is the exam's own level: under e-LOND and spending it is the level $\alpha_t$ fixed when the exam starts, and the bet is tuned to it. The deployed gate uses instead the level-tuned grid mixture of Theorem~\ref{thm:adaptive}, a mixture of such processes, which Proposition~\ref{prop:valid} also covers. The exam admits at the first $n$ with $K_n\ge1/\alpha$. For futility it runs a second, downward bet $K^{\mathrm f}_n=\prod_i(1-\lambda^{\mathrm f}_i(X_i-m_{\mathrm f}))$ against $m_{\mathrm f}=(1+\delta_{\mathrm f})/2$ with $\delta_{\mathrm f}\ge\delta$, and rejects when $K^{\mathrm f}_n\ge1/\alpha_{\mathrm f}$; otherwise it rejects at $N_{\max}$.

\begin{proposition}[Validity]\label{prop:valid}
If $\mathbb{E}[X_i\mid\mathcal F_{i-1}]\le m_0$ for every $i$, then for any stopping rule, including the futility rule and any futility margin $\delta_{\mathrm f}$, $\Pr(\text{admit})\le\Pr(\sup_n K_n\ge1/\alpha)\le\alpha$, and the stopped wealth $K_\tau$ is an e-value: $\mathbb{E}[K_\tau]\le1$.
\end{proposition}
The proof (Appendix~\ref{app:proofs}) is Ville's inequality for the nonnegative supermartingale $K_n$. Two remarks matter in practice. First, the condition is on conditional means, not on identically distributed variants, so drift in a provider's model during an exam does not break validity as long as the arms run concurrently and the skill's conditional lift stays at or below $\delta$; a comparator taken from an earlier run does break it (Section~\ref{sec:res-validity}). Second, the futility margin is free in type-I error: it only removes admissions. Testing futility at $\delta_{\mathrm f}=\delta$, the obvious default, almost never fires when $\Delta=\delta$ exactly, so useless skills run to $N_{\max}$; we use $\delta_{\mathrm f}=\PFutMFrac$ (\PFutM{} percentage points), half the design lift. The simpler cap $\lambda_i\le m_0$ is close to $c=1/2$; we report it as a separate variant.

For pass/fail verifiers with $r=1$ and $\delta=0$, Wald's SPRT on discordant pairs tests $p=\Pr(D_i=1\mid D_i\ne0)\le1/2$; its likelihood ratio is also a test supermartingale, so it satisfies Proposition~\ref{prop:valid} and gives an e-value. It cannot express a margin $\delta>0$ or graded scores, and it needs a pre-specified alternative.

\subsection{Library and resubmissions (R2, R3)}
A natural design stops each exam at $K_n\ge1/\alpha$ and then applies e-BH \citep{wang2022ebh} to the $m$ exams of an admission window. This fails.
\begin{remark}[Stopped e-values are too small for e-BH]\label{prop:ebh}
If each exam stops at the first $n$ with $K_n\ge1/\alpha$ and every bet satisfies $\lambda_i\le\bar\lambda<1/m_0$, then $K_\tau<(1+\bar\lambda(1-m_0))/\alpha$. e-BH at level $\alpha$ over $m$ such exams admits nothing unless more than $m/(1+\bar\lambda(1-m_0))$ of them crossed. At $\delta=0$ this is $80\%$ of the window for PrPl with $c=1/2$ ($\bar\lambda=1/2$), $73\%$ for $c=3/4$, and $67\%$ for aGRAPA with $c=1/2$ ($\bar\lambda=c/m_0=1$). Every admissible bet has $\bar\lambda\le1/m_0$, and $K_\tau<1/(m_0\alpha)$ holds even at equality, so at least a fraction $m_0$ of the window must cross, half of it at $\delta=0$.
\end{remark}
\citet{wang2022ebh} already observe that the all-or-nothing e-value $\mathbf 1\{p\le\alpha\}/\alpha$ carries little information for later analyses, unlike a generic e-value; the remark makes explicit what this means for e-BH over exams stopped at $1/\alpha$. Valid alternatives either stop at the strictest e-BH level, $m/\alpha$, or run every exam to $N_{\max}$ before applying e-BH; both cost variants. Stopped e-BH \citep{wang2025stopped} and e-BH on suprema of e-processes \citep{tavyrikov2025carefree} are principled versions of these fixes. Exams in a registry arrive over time, which suits online control. e-LOND \citep{xu2024online} assigns exam $t$ the level $\alpha_t=\alpha\gamma_t(R_{t-1}+1)$, where $R_{t-1}$ counts earlier admissions and $\sum_t\gamma_t\le1$; the exam then stops at $K_n\ge1/\alpha_t$. FDR stays below $\alpha$ under arbitrary dependence among valid e-values; here that means any dependence between exams that leaves each null exam's conditional null intact given the registry's history when it starts. The threshold is known when the exam starts, so it composes with sequential stopping. Using only admissions known at the start of a window keeps $\alpha_t$ predictable.

\paragraph{The unit of error.} e-LOND counts exams and admissions; R2 is stated for admitted skills. They agree under one convention, which the gate and all our simulations use. An exam tests one version, for one agent, against one comparator, and ends with at most one admission; every commitment consumes an e-LOND index and an attempt index of its lineage, whether or not the exam finishes, so an author cannot abort and restart at the same level. A lineage is admitted at most once against a given comparator: once a version is admitted, later versions of the same lineage are replacements, tested against the admitted version as comparator $s_0$, and their exams continue the lineage's attempt count. A retry of the same version is a new exam. Exams that overlap in time are separate hypotheses whose levels interact only through e-LOND. Each admission therefore lists exactly one (version, agent, comparator) triple, and the FDR over exams is the FDR over admitted triples. A skill admitted for several agents counts once per agent; a registry that wants FDR over distinct skills must count admissions by skill in $R_{t-1}$ and test each skill once per comparator. For example, suppose version~1 of a skill is useless and stops for futility, version~2 is helpful and is admitted, and version~3 is an edit that adds nothing. Version~3 is tested against version~2; its null is true, so admitting it would be a false discovery, counted once in the FDR and once against the lineage's budget. E22 checks that the gate, the proof and the simulations count the same unit (Appendix~\ref{app:more}).

FDR control does not stop a determined author: each resubmission of a useless skill is a fresh exam with a fresh level. We therefore spend alpha over a skill's lineage.
\begin{proposition}[Lineage alpha-spending]\label{prop:lineage}
If the $k$-th exam of a lineage runs at level $\alpha_k^{\text{lin}}=\min(\alpha_t,\ \alpha\gamma^{\text{lin}}_k)$, fixed when it starts, with $\sum_k\gamma^{\text{lin}}_k\le1$, then the probability that any version with $\Delta\le\delta$ is admitted is at most $\alpha$, whatever helpful versions the lineage also contains and however the author chooses versions from earlier results, and the e-LOND FDR guarantee still holds.
\end{proposition}
The rule needs a lineage definition: the registry links resubmissions by author key, task class and skill similarity. Linking is a policy choice outside the statistics; an author who evades linking with new identities is a sybil problem that stake or identity verification must address.

\subsection{Dual-arm credential and verification (R4)}
Each arm's transcripts are hashed into a Merkle tree, and a signed credential binds the skill digest, the commitment and reveal, $\theta$, both roots and the statistics $(N,\hat\Delta,\mathrm{LCB},\text{decision})$ (Appendix~\ref{app:schema}). Algorithm~\ref{alg:verify} (Appendix~\ref{app:schema}) is what a verifier runs. It re-executes only the deterministic verifiers on recorded final states, not the agent, so it is cheap (Section~\ref{sec:res-leak}). It cannot show that the recorded behaviour came from the stated model; that needs attested capture, which we treat as an extension.



\paragraph{What the credential should report.} A credential issued at the stopping time reports the estimate $\hat\Delta=\bar D_N$, which is biased upward because exams stop when the estimate is high, and the anytime-valid lower bound obtained by inverting Eq.~\eqref{eq:wealth} over $m_0$, which keeps its stated time-uniform coverage under the conditions of Proposition~\ref{prop:valid} but sits just above $\delta$ at admission (Section~\ref{sec:res-report}). The credential should lead with the bound; a registry that wants an informative effect size can continue the exam after admission at no cost to validity.
